\BourbakiFrontMatter{引言}

本章致力于研究环的某些类以及这些环上的模。这一研究背后有若干主题，例如分类或分解问题，以及对子对象和态射集的描述。

我们实际上只在合理的有限性假设下处理这些问题，这正是本章从模的概念以及诺特环和阿廷环的概念开始的原因。

就环而言，我们给出若干使人能够理解阿廷环的结果：

\begin{enumerate}
\def\labelenumi{\alph{enumi})}
\item
  阿廷环的雅各布森根是它的幂零根（\(\S10\)），而相应的商环是半单环（\(\S8\)）。
\item
  半单环同构于有限族单环的乘积（§7）。
\item
  由韦德伯恩定理（VIII, p.~120, 定理 1），单环同构于体上的一个矩阵代数。
\end{enumerate}

在森田等价（\(\S6\)）的意义下，是单环的有限次数代数由它在一个称为布饶尔群的交换群中的类所确定。我们给出这个群的若干描述（\(\S15\) 和 \(\S16\)），以及一些例子（\(\S18\) 和 \(\S19\)）。

对于模，我们给出两个自然的分解：第一个借助合成列（I, §4, No.~7, p.~41, 定义 9）表述，由若尔当--赫尔德定理（I, §4, No.~7, p.~43, 定理 6）提供。第二个对应于直和，在有限长度模的情形由克鲁尔--雷马克--施密特定理（VIII, p.~37, 定理 2）给出，在这里我们由阿祖马亚关于半原始模的一个结果（VIII, p.~32，

定理 1）推出它。我们还考虑与模相伴、对上述分解具有加性表现的不变量；§11 中描述的格罗滕迪克群是这些不变量的泛问题的解。在研究环上模的结构时，环同构的概念有利地被森田等价所取代。

对于半单模，即单模的直和，模的描述这一概念（VIII, p.~69, 定义 5）使我们既能描述从该模出发的同态，又能描述它的子模。

作为说明，在 §21 中我们考虑有限群的代数的情形，其模对应于群的线性表示。

卷末附有一篇从上一版照录的历史笔记，它追溯了此处发展的许多概念是如何产生的。

\BourbakiMainMatter
\setcounter{chapter}{7}
